\subsection{Symplectic vector bundles}

\bdefn

    A {\emphcolor symplectic vector bundle} over a manifold $M$ is a pair $(E,\omega)$ where $\pi\colon E\to M$ is a real vector bundle and $\omega\in\Gamma(\Lambda^2(E^*))$
    is a smooth tensor field such that for all $p\in M$, $\omega_p$ is a symplectic form on $E_q=\pi^{-1}q$.

    Two symplectic vector bundles $(E_i,\omega_i)$, for $i=1,2$, are isomorphic if there is a vector bundle isomorphism $\Psi\colon E_1\to E_2$ such that $\Psi^*\omega_2=\omega_1$.

\edefn

\bexam

    If $E\to M$ is a vector bundle, $E\oplus E^*$ has a canonical structure as a symplectic vector bundle via
    $$ \Omega_\can(v\oplus\phi,u\oplus\psi) = \psi(v) - \phi(u) . $$

\eexam

\blemm

    If $(E,\omega)$ is a symplectic vector bundle over $M$, $F\subseteq E$ a vector subbundle, then $F^\omega\subseteq E$ is also a vector subbundle defined pointwise by
    $$ (F^\omega)_p = (F_p)^{\omega_p} . $$

\elemm

\bproof

    Define
    $$ \hat\omega\colon E \to F^*,\qquad \hat\omega(\epsilon)(\eta) = \omega(\epsilon,\eta)\quad\hbox{for $\epsilon\in E,\eta\in F$} , $$
    then $\hat\omega$ is a smooth vector bundle morphism and $F^\omega=\ker\hat\omega$, which is a subbundle iff $F$ has constant rank.
    And indeed it does, its rank at each point is $\dim(F_p)^\omega=\dim E_p-\dim F_p=\rank E-\rank F$.
    \qed

\eproof

\bprop

    Let $(E,\omega)$ be a symplectic vector bundle of rank $2n$ over $M$, then:
    \benum
        \item $E$ has a system of local trivializations whose transition maps take values in $\Sp(2n)$.
        \item Every point $p\in E$ has a neighborhood $U\subseteq M$ such that $(\pi^{-1}U,\omega)\subseteq(E,\omega)$ is isomorphic to the symplectic bundle $(U\times\bR^{2n},\omega_0)$.
        \item Every point $p\in E$ has a local frame which forms a symplectic basis at each point (called a {\emphcolor symplectic frame}).
    \eenum
    Furthermore (2) and (3) are equivalent to $(E,\omega)$ being symplectic, and (1) gives $E$ a symplectic structure.

\eprop

\bproof

    Clearly (2) and (3) are equivalent, as a symplectic frame $\epsilon_1,\dots,\epsilon_n,\delta_1,\dots,\delta_n$ over $U$ defines an isomorphism from $(U\times\bR^{2n},\omega_0)$ by
    $$ (p,(x_1,\dots,x_n,y_1,\dots,y_n)) \mapsto \sum_{i=1}^n x_i\epsilon_i(p) + y_i\delta_i(p) , $$
    and conversely such an isomorphism induces a symplectic frame by considering the image of the symplectic frame $x_1,\dots,x_n,y_1,\dots,y_n$ of $U\times\bR^{2n}$.
    Clearly (2) implies that $(E,\omega)$ is symplectic, as through these isomorphisms we see that it is pointwise antisymmetric, bilinear, and nondegenerate.

    We show that $(E,\omega)$ being symplectic implies (3).
    Let $\epsilon_1,\delta_1$ be local sections around $p\in E$ such that $\omega(\epsilon_1,\delta_1)\neq0$.
    We can find them by choosing $e_1,f_1\in E_p$ such that $\omega_p(e_1,f_1)\neq0$ and extending them to local sections, and then by continuity there is a neighborhood where $\omega(\epsilon_1,\delta_1)\neq0$.
    Then $\epsilon_1,\delta_1$ generate a symplectic subbundle of $E|_U$, and taking its symplectic complement we obtain a symplectic subbundle of rank $2n-2$.
    By induction we get a local symplectic frame.

    Now we show (3) implies (1).
    If $\set{\epsilon^\alpha_i,\delta^\alpha_i}_{i\leq n}$ and $\set{\epsilon^\beta_i,\delta^\beta_i}_{i\leq n}$ are local symplectic frames over $U_\alpha,U_\beta$ respectively,
    then the transition maps at each point corresponds to a transition map between symplectic bases.
    This is a symplectic matrix, and so the transition map takes values in $\Sp(2n)$ as desired.

    Now assume that $(U_\alpha,\Phi_\alpha)_\alpha$ is a system of local trivializations whose transition maps take values in $\Sp(2n)$, then they induce a symplectic structure on $E$.
    Indeed, taking an $\alpha$, let $\epsilon_i^\alpha=\Phi_\alpha^{-1}\circ x_i$ and $\delta_i^\alpha=\Phi_\alpha^{-1}\circ y_i$ be the pullbacks of the symplectic frames of $U_\alpha\times\bR^{2n}$.
    Then define $\omega_\alpha$ on $\pi^{-1}U_\alpha$ via this local frame.
    Because the transition maps take values in $\Sp(2n)$, this means that $\omega_\alpha$ and $\omega_\beta$ agree on their overlaps, and thus glue together to form a symplectic structure on $E$.
    \qed

\eproof

\bdefn

    A {\emphcolor complex structure} on a vector bundle $E\to M$ is an automorphism $J$ of $E$ such that $J^2=-\id$.
    If $(M,\omega)$ is symplectic, $J$ is said to be {\emphcolor compatible} with $\omega$ if $J_p$ is compatible with $\omega_p$ for all $p\in M$.
    Any compatible pair $(\omega,J)$ defines a metric $g_J$ on $E$ bt
    $$ g_J(v,u) = \omega(v,Ju) . $$
    Such a triple $(\omega,J,g)$ is called a {\emphcolor Hermitian structure} on $E$.

\edefn

\bprop[name={homotopsympbund}]

    Let $E\to M$ be a rank $2n$ vector bundle.
    \benum
        \item For every symplectic structure $\omega$ on $E$, the space of $\omega$-compatible and $\omega$-tame complex structures, denoted $\J(E,\omega)$ and $\J_\tau(E,\omega)$ respectively,
        are nonempty and contractible.
        \item Let $J$ be a complex structure on $E$, then the space of symplectic structures on $E$ compatible with $J$ (or which tame $J$) is nonempty and contractible.
        \item Let $g$ be a metric on $E$, then
        $$ \S(E) \xtof{\omega\mapsto J_{g,\omega}}[J\mapsto\omega_{g,J}] \J(E) $$
        form a homotopy equivalence.
    \eenum

\eprop

\bproof

    For (1) consider the maps
    $$ \J(E,\omega) \xtof{J\mapsto g_J}[g\mapsto J_{g,\omega}] \Met(E) . $$
    These are homotopy inverses, by \refmath[corollary]{compatlincompcontract}.
    And $\Met(E)\subseteq\Sym^2(V^*)$ is a convex subset of a vector space, and thus contractible.

    Now we need to show that $\J_\tau(E,\omega)$ is contractible.
    Let $J_0\in\J(E,\omega)$, and denote by $\iprod{\bul,\bul}_0=g_{J_0}$ its associated metric on $E$; and denote its associated norm by $\norm\bul_0$.
    This induces the operator norm $\endo(E_p)$ for $p\in M$, which we denote by $\norm\bul_0$ as well.
    Now define
    $$ \A = \set{A\in\endo(E)}[AJ_0+J_0A=0,\ \norm{A_p}_0<1\hbox{ for all $p\in M$}] . $$
    Then define
    $$ \F\colon\J_\tau(E,\omega) \to \A,\qquad \F(J) = (\id+J_0J)(\id-J_0J)^{-1} , $$
    notice that $\id-J_0J$ is invertible because if $J_0Jv=v$ then $Jv=-J_0v$, but then
    $$ \omega(v,Jv) = -\omega(v,J_0v) < 0 , $$
    contradicting tameness.
    It is then (non-trivial) linear algebra to show that this is well-defined.

    $\F$ is a homeomorphism with inverse
    $$ \F^{-1}\colon\A\to\J_\tau(E,\omega),\qquad \F^{-1}(A) = J_0(\id+A)^{-1}(\id-A) . $$
    Thus $\J_\tau(E,\omega)$ is homeomorphic to the convex $\A$, so it is contractible.

    In (2), the existence of a symplectic structure compatible with $J$ is equivalent to the existence of a metric $g$ with respect to which $J$ is skew-adjoint.
    Take any metric $g$ and consider $\frac12(g+J^*g)$, giving the desired result.
    Contractibility follows from the fact that the space in question is convex.

    (3) follows from \refmath[proposition]{homotopymapscompsymp}; the same definitions form the desired homotopy inverses.
    \qed

\eproof

\blemm

    Suppose $J_t$ is a smooth family of complex structures on $E$.
    Then there is a smooth family of isomorphisms $\Phi_t\colon E\to E$ such that $\Phi_t^*J_t=J_0$.

\elemm

\bproof

    Define
    $$ B_t = \frac12\partial_tJ_tJ_t , $$
    and then consider
    $$ [B_t,J_t] = B_tJ_t - J_tB_t = \frac12\partial_tJ_tJ_t^2 - \frac12J_t\partial_tJ_tJ_t = -\frac12\partial_tJ_t - \frac12\partial_tJ_t = \partial_tJ_t , $$
    here we use the fact that $J_t\partial_tJ_tJ_t=\partial_tJ_t$, this is because differentiating $J_t^2=-\id$ gives
    $$ \partial_tJ_tJ_t + J_t\partial_tJ_t = 0 \implies J_t\partial_tJ_tJ_t = \partial_tJ_t . $$

    Then consider the linear ODE
    $$ \partial_t\Phi_t = B_t\Phi_t,\qquad \Phi_0 = \id , $$
    which has a unique smooth solution.
    Furthermore $\Phi_t$ is an isomorphism (this is a classical result in ODE theory).

    Now define $D_t=J_t\Phi_t-\Phi_tJ_0$, then
    $$ \partial_tD_t = \partial_tJ_t\Phi_t + J_tB_t\Phi_t - B_t\Phi_tJ_0 = (\partial_tJ_t + J_tB_t - B_tJ_t)\Phi_t + B_tD_t = B_tD_t . $$
    Thus $D_t$ solves the ODE $\partial_tD_t=B_tD_t$ with the initial value $D_0=0$.
    Because $D_t=0$ solves this problem, by uniqueness $D_t=0$ as desired.
    \qed

\eproof

\bthrm

    Let $(E_i,\omega_i)$ be symplectic vector bundles over $M$, and let $J_i$ be $\omega_i$-tame for $i=1,2$.
    Then $(E_1,\omega_1)$ and $(E_2,\omega_2)$ are isomorphic iff the complex vector bundles $(E_1,J_1)$ and $(E_2,J_2)$ are isomorphic.

\ethrm

\bproof

    Suppose there is an isomorphism $\Psi\colon E_1\to E_2$ with $\omega_1=\Psi^*\omega_2$.
    Then both $J_1$ and $\Psi^*J_2$ are $\omega_1$-tame, so by (1) in \refmath[proposition]{homotopsympbund}, there is a smooth family of complex structures $J_t\colon E_1\to E_2$ which are $\omega_1$-tame
    and $J_0=\Psi^*J_2$.
    By the above lemma there is a smooth family of isomorphisms $\Phi_t\colon E_1\to E_1$ such that $\Phi_t^*J_t=J_1$.
    Then $\Psi\circ\Phi_0\colon E_1\to E_2$ is an isomorphism and
    $$ (\Psi\circ\Phi_0)^*J_2 = \Phi_0^*\Psi^*J_2 = \Phi_0^*J_0 = J_1 , $$
    as desired.

    Conversely suppose there is a bundle isomorphism $\Psi\colon E_1\to E_2$ such that $\Psi^*J_2=J_1$.
    Then $\Psi^*\omega_2,\omega_1$ are symplectic forms on $E_1$ which tame $J_1$, and therefore so is
    $$ \omega_t = t\omega_1 + (1-t)\Psi^*\omega_2 , $$
    and so $\Psi^*\omega_2,\omega_1$ are connected by a smooth family of symplectic forms.
    By \refmath[proposition]{sympfamfromiso} there is a smooth family of isomorphisms $\Phi_t\colon E_1\to E_1$ such that $\Phi_1=\id$ and $\Phi_t^*\omega_t=\omega_1$.
    Thus $\Phi_0^*\Psi^*\omega_2=\omega_1$, meaning $\Psi\circ\Phi_0$ is the desired isomorphism.
    \qed

\eproof

